% © 2026 Ing. David Jaroš — CC BY-NC-ND 4.0
% Canonical biquaternionic curvature.
\section{Biquaternionic curvature}
\label{sec:canonical:biquaternion_curvature}

For the general two-sided biquaternionic derivative
\begin{equation}
 D_\mu\Theta=\partial_\mu\Theta+A_\mu\Theta-\Theta B_\mu,
\end{equation}
define
\begin{align}
 F^A_{\mu\nu}
 &=\partial_\mu A_\nu-\partial_\nu A_\mu+[A_\mu,A_\nu],\\
 F^B_{\mu\nu}
 &=\partial_\mu B_\nu-\partial_\nu B_\mu+[B_\mu,B_\nu].
\end{align}
Then
\begin{equation}
 \boxed{
 [D_\mu,D_\nu]\Theta
 =F^A_{\mu\nu}\Theta-\Theta F^B_{\mu\nu}.}
 \label{eq:canonical:bimodule-curvature}
\end{equation}
The one-sided formula is recovered only as the special case $B_\mu=0$.

In the torsion-free compatible branch, antisymmetric integrability requires
\begin{equation}
 F^A_{\mu\nu}\Theta=\Theta F^B_{\mu\nu}.
\end{equation}
For invertible $\Theta$, the curvatures may be nonzero but must be conjugate:
\begin{equation}
 F^A_{\mu\nu}=\Theta F^B_{\mu\nu}\Theta^{-1}.
\end{equation}
This is why the two-sided action avoids the generic flatness obstruction of a
naive one-sided regular representation.

The Lorentz frame curvature built from the uniquely reconstructed
$\omega=\mathring\omega(e)+K(T)$ represents the corresponding Riemann--Cartan
geometry.  For $T=0$ it is the usual Riemann curvature of the Levi--Civita
connection.  The pure Lorentz identification is exactly
$A_\mu=\Omega_\mu$, $B_\mu=-\Omega_\mu^\ddagger$ modulo a cancelling
central term, but that branch is a concurrent-vector no-go for generic curved
GR.  What remains open is the derivation and algebraic elimination of a
nontrivial relative bimodule component from the UBT action.

The flat special-relativistic branch has zero curvature and admits an inertial
gauge with $A_\mu=B_\mu=0$, hence $D_\mu=\partial_\mu$.
